%%%%%%%%%%%%%%%%%%%%%%%%%%%%%%%%%%%%%%%%%%%%%%%%%%%%%%%%%%%%%%%%%
\chapter{Unit 7: Large Language Models in Biomedical Research}
\label{chap:unit7}
%%%%%%%%%%%%%%%%%%%%%%%%%%%%%%%%%%%%%%%%%%%%%%%%%%%%%%%%%%%%%%%%%

\textbf{Total Time: 5 hours}

This unit prepares participants to use large language model (LLM)
based AI agents as rigorous, accountable tools in biomedical data
science research. The unit is \emph{load-bearing}: its deliverables
are components of the 2026 Data Challenge final report. Students
work with the same NCHS vital statistics database used throughout
the bootcamp, using an LLM ensemble to assist with data navigation,
methodological decision-making, grant review, and manuscript
preparation, while maintaining human accountability for every output.

Three principles organize the unit. First, invalidation
(a framework broader than ``hallucination'' that encompasses
factual, logical, normative, and structural breaches) is
mathematically inevitable in any LLM-generated output; human
verification is therefore not optional. Second, cross-agent
critique reduces invalidation rates below what any single model
achieves, but introduces its own ethical tradeoffs that must be
actively managed. Third, race and ethnicity in biomedical data are
social constructs whose reification in data systems has real
consequences; LLMs applied to historical vital statistics data are
particularly prone to normative invalidation when treating racial
categories as stable natural kinds.

%================================================================
\section{The Ethics of AI Agents}
\label{sec:7.1}
%================================================================

\textbf{Time: 1 hour (Instructional: 50 minutes, Protocol Construction: 10 minutes)}

This lesson examines the ethical dimensions of using LLM-based AI
agents in biomedical data science. Unlike single-turn chatbot
interactions, agentic workflows involve planning, tool use,
persistent memory, and multi-agent coordination, substantially
expanding the ethical surface area of AI-assisted research.
Students develop a framework for identifying and classifying
invalidation risks, analyze the specific ethical challenges posed
by historical vital statistics data with inconsistently encoded
racial and ethnic categories, and produce a \emph{Research AI
Governance Protocol} that governs their own use of LLM tools
throughout Sections~\ref{sec:7.2} and~\ref{sec:7.3}.

The lesson proceeds in four instructional blocks:

\begin{enumerate}
    \item \textbf{Agent architecture and expanded ethical surface
    (15 min).} What distinguishes an AI agent from a chatbot:
    planning, tool use, memory, and multi-agent coordination.
    Each capability adds a distinct ethical surface. The discursive
    network framework explains why agent-generated statements that
    enter manuscripts acquire authority independent of their accuracy,
    making verification a structural responsibility rather than an
    optional quality check.

    \item \textbf{Invalidation: a taxonomy broader than hallucination
    (15 min).} The four types of invalidation (factual, logical,
    normative, and structural) are introduced using NCHS-specific
    examples. Factual: an LLM names a variable not present in the
    DVS schema as a named field. Logical: an LLM recommends a rate
    denominator inconsistent with the research question's population
    definition. Normative: an LLM interprets racial disparities in
    birthweight as reflecting biological differences between racial
    groups rather than the consequences of racism. Structural: an
    LLM queries APGAR score columns for records before 1977, when
    those columns did not exist in the flat file. Each type requires
    a distinct verification strategy.

    \item \textbf{Race-consciousness in biomedical data science
    (15 min).} Race is a social construct with no valid biological
    basis, but its reification in law, medicine, and data systems
    makes it real in its consequences. Birth outcomes differ by
    recorded race not because of intrinsic racial biology but because
    of racism: differential access to prenatal care, chronic stress
    from discrimination, environmental exposures in historically
    redlined neighborhoods, and differential clinical treatment.

    The specific NCHS data students are working with encodes race
    through a system that changed across the study period. From 1969
    through 1977, race was assigned from a fixed set (White, Black,
    American Indian, Chinese, Japanese, Other), and for parents
    of different recorded races the child's race followed the
    mother's. Assignment rules were revised in 1978, inconsistently
    across states. The ``Hispanic'' ethnic category appeared late
    in the period and was collapsed into ``White'' in some years.
    The \texttt{child\_race\_recode\_3} field changed its encoding
    midway through the study period. Any LLM query about race in
    this dataset risks producing normative invalidation by treating
    these categories as stable, comparable, and biologically grounded
    across the full 1969--1986 span.

    \item \textbf{Ethics checklist and governance protocol
    construction (20 min).} The seven risk categories from the
    ethics framework (epistemic, accountability and authorship,
    provenance and reproducibility, confidentiality, bias and
    fairness, security, and sustainability) are reviewed with
    emphasis on the tradeoffs specific to this bootcamp context.
    Multi-agent critique tradeoffs: privacy amplification, epistemic
    homogenization, compute burden, and attribution diffusion.
    Students construct the Research AI Governance Protocol that
    governs their subsequent work.
\end{enumerate}

\subsection{Learning Objectives}

\begin{enumerate}
    \item Define an AI agent in terms of planning, tool use,
    memory, and multi-agent interaction, and explain how each
    capability expands the ethical surface area relative to
    single-turn LLM use.

    \item Apply the four-type invalidation taxonomy (factual,
    logical, normative, and structural) to concrete scenarios
    involving the NCHS vital statistics database, and identify
    the appropriate verification strategy for each type.

    \item Articulate why race and ethnicity in historical US
    vital statistics data are social constructs whose reification
    has real health consequences, and explain the specific coding
    inconsistencies in the NCHS 1969--1986 dataset that make
    naive racial comparisons normatively invalid.

    \item Construct a Research AI Governance Protocol specifying
    provenance logging, invalidation classification, race-variable
    handling, escalation thresholds, and a compliant disclosure
    statement for the team's subsequent AI-assisted work.
\end{enumerate}

\subsection{Assessment Instrument}

\textbf{Item 1 --- Invalidation Taxonomy Application (in-session,
8 minutes):}
The instructor presents four NCHS-specific LLM output scenarios,
one per invalidation type. For each scenario, teams identify:
(a) the invalidation type, (b) why it constitutes that type and
not another, and (c) the verification method that would detect
and correct it. Scenarios are drawn from realistic ensemble
outputs on the vital statistics database.

\textbf{Item 2 --- Research AI Governance Protocol
(in-session construction, 10 minutes; referenced throughout
Sections~\ref{sec:7.2} and \ref{sec:7.3}):}
Each team produces a one-page governance protocol containing
all five required components:
\begin{enumerate}
    \item \textbf{Provenance logging convention:} how prompts,
    model identifiers, and timestamps will be recorded for every
    ensemble interaction.

    \item \textbf{Invalidation log format:} the template for
    recording invalidation type, description, and resolution for
    each detected LLM failure.

    \item \textbf{Race variable handling clause:} specific
    commitments about how race and ethnicity fields will be
    treated in LLM queries given the documented coding
    inconsistencies of the NCHS dataset. At minimum: LLM queries
    involving race variables must inject the relevant coding
    history as context; temporal comparisons of racial category
    rates require explicit justification; normative claims about
    racial disparities must be flagged for human review.

    \item \textbf{Escalation threshold:} the conditions under
    which LLM output will not be used without independent
    verification or correction (e.g., any factual claim about
    NCHS variable definitions; any normative statement about
    population health disparities; any generated code before
    execution).

    \item \textbf{Draft disclosure statement:} a single sentence
    in the form specified by the ethics framework, naming tools
    and tasks, to be finalized at the end of Section~\ref{sec:7.3}.
\end{enumerate}

The governance protocol is a live document. It is updated
throughout Sections~\ref{sec:7.2} and~\ref{sec:7.3} as
invalidations are encountered and resolved, and its disclosure
statement is finalized as the unit's final deliverable.

%================================================================
\section{AI Agents for Technical Tasks}
\label{sec:7.2}
%================================================================

\textbf{Time: 2 hours (Instructional: 40 minutes, Activity Work:
80 minutes)}

This lesson teaches students to use an LLM ensemble as a rigorous
technical collaborator on analytical tasks grounded in the NCHS
vital statistics database. The session has four sequential blocks.
Block~A derives the mathematical basis for multi-agent consensus
and produces the agent-count commitments that govern subsequent
activities. Blocks~B, C, and~D form an analytical chain: outputs
from each block feed the next, and their collective products are
load-bearing components of the final data challenge report.

Students enter this session with a completed model report from
Unit~4, which documents a predictive model for the data
challenge outcome (underweight birth rate or infant mortality
rate by Texas county) including the feature selection memo,
model comparison table, residual diagnostics, and limitations
section. This report is the primary document the LLM ensemble
will be asked to interrogate. The governance protocol
from Section~\ref{sec:7.1} is active from the start of this
session; every ensemble interaction concerning the Unit~4
model is logged in the invalidation log.

The governing framework is the \emph{Flaws-of-Others} (FOO)
algorithm, in which multiple agents independently produce outputs,
each critiques the others' responses, and a harmonizer synthesizes
the critiques into a refined result. The minimum number of
mutually detecting agents required to reduce the long-run
false-statement share below a tolerance $\varepsilon$ is:
%
\begin{equation}
    n_{\min} =
    \left\lceil
        1 + \frac{p\!\left(\dfrac{1}{\varepsilon} - 1\right) - q}{d}
    \right\rceil
    \label{eq:nmin}
\end{equation}
%
where $p$ is the per-agent invalidation rate, $q$ is the internal
self-repair rate ($q = 0.05$ reference value), and $d$ is the
cross-detection probability ($d = 0.19$ reference value). Students
use this formula to make defensible, quantitative decisions about
ensemble size for each activity.

\subsection*{Block A --- FOO Framework: Derivation and Calibration
(20 minutes instructional)}

The derivation proceeds from first principles. A single agent has
an effective invalidation rate $p$ (combining intrinsic corruption
and fabrication). Internal self-repair operates at rate $q$.
Each additional cross-checking agent adds detection hazard $d$
to the effective correction rate. The steady-state false-statement
share with $n$ mutually detecting agents is:
%
\begin{equation}
    \pi_F^{(n)} = \frac{p}{p + q + (n-1)d}
    \label{eq:steadystate}
\end{equation}
%
Setting $\pi_F^{(n)} \leq \varepsilon$ and solving for $n$ yields
Eq.~(\ref{eq:nmin}). The invalidation floor (proved formally in the
underlying theory) establishes that some error rate is
mathematically inevitable in any finite-loss LLM; the FOO
architecture is the engineering response to this impossibility
result.

Students compute $n_{\min}$ for three task categories they will
encounter in this session using the reference values above:
{\footnotesize
\begin{center}
\begin{tabular}{llccc}
\toprule
Task category & Example & $p$ & $\varepsilon$ & $n_{\min}$ \\
\midrule
Factual schema lookup & Identify correct byte range for birthweight & 0.10 & 0.10 & \emph{compute} \\
Methodological judgment & Choose projection model
    for county-level rates & 0.30 & 0.05 & \emph{compute} \\
Code generation & SQL query for core
    research question & 0.20 & 0.05 & \emph{compute} \\
\bottomrule
\end{tabular}
\end{center}
}
%
The computed values are entered into the governance protocol as
the team's agent-count commitments for Blocks~B, C, and~D.

\subsection*{Block B --- The Texas Births Problem (35 minutes)}

Before the bootcamp, participants were asked to count live births
in Texas in 1969 to mothers residing in Texas. Teams compare their
answers. Historically, different correct implementations of this
query return different counts. The task is to diagnose the source
of disagreement using the LLM ensemble.

Teams query the ensemble with the NCHS data dictionary and ask it
to: (a)~identify every field relevant to the filter, (b)~flag
any ambiguities in the dictionary language, and (c)~write a query
implementing the most defensible interpretation. Teams execute
the generated query and compare counts across the class.

%The instructor then surfaces the documented sources of ambiguity:
% the \texttt{resident} flag records whether the mother is a resident
% of the \emph{reporting area} (not necessarily Texas); the
% place-of-residence field encodes the mother's state; and the
% reporting area field records where the birth was registered. A
% count restricted to records where the birth \emph{occurred} in
% Texas differs from one restricted to mothers \emph{residing} in
% Texas, and both differ from one using the resident flag. The
% ``live birth'' filter has no explicit flag; it is the absence
% of certain mortality indicators that were encoded differently
% across years.

The deliverable is annotated, executable code (Python or SQL) that
implements the correct filter with comments explaining each
disambiguation decision. This code is a candidate component of
the reproducible pipeline developed in Unit~5.

\textbf{Model report critique (15 min).} Each team submits
their Unit~4 model report to the ensemble and prompts it to
identify any claim in the methods section that is (a)~not
reproducible from the written description alone, (b)~in
tension with a stated limitation, or (c)~likely to require
verification against the NCHS data dictionary. Teams record
each flagged claim in the invalidation log with an invalidation
type classification from Section~\ref{sec:7.1}. For each
factual or structural invalidation identified, teams verify
the ensemble's critique against the Unit~4 model code before
accepting or rejecting it. Teams that find no invalidations
in their Unit~4 report should examine whether the ensemble
is producing false-negative outputs (a structural invalidation
of the ensemble itself) and document their reasoning.

\subsection*{Block C --- Race Variable Temporal Inconsistency
(15 minutes)}

Building directly on Section~\ref{sec:7.1}, students encounter
the race variable problem in the actual data. Teams query the
ensemble with the race coding history injected as context and ask
it to: (a)~identify which race-related variable or combination
is most defensible for the projection analysis, (b)~specify which
temporal comparisons are valid and which are not, and
(c)~identify any normative claims in its own response.

The key learning outcomes are: (1)~students observe whether the
ensemble correctly identifies the temporal inconsistency and flags
its own normative claims, (2)~students discover how injecting the
coding history as context changes ensemble output, demonstrating
retrieval-augmented generation as a concrete mitigation for
normative and structural invalidation, and (3)~students update the
race variable handling clause in their governance protocol based
on what they actually encountered.

\subsection*{Block D --- Methodological Decision Support
(60 minutes)}

Teams use the ensemble to interrogate the model family choice
documented in their Unit~4 report. The core prompt pattern
is: present the ensemble with the Unit~4 feature set, the
outcome distribution description, and the model comparison
table, then ask it to (a)~identify which alternative model
from the comparison was most appropriate given the stated
distributional properties of the outcome, and (b)~identify
any model family that was not evaluated in the Unit~4
comparison but would be worth considering given the data
structure. Teams apply Eq.~(\ref{eq:nmin}) to determine
the appropriate ensemble size for this task. Every ensemble
recommendation is logged and verified against the Unit~4
model diagnostics before being accepted. If the ensemble
recommends a model family that the student already
evaluated and rejected in Unit~4, teams document whether
the ensemble's reasoning accounts for the diagnostic
evidence that motivated the rejection.

\subsection{Learning Objectives}

\begin{enumerate}
    \item Derive the FOO minimum-agent formula from the steady-state
    false-share equation and compute $n_{\min}$ for tasks of
    varying epistemic difficulty, applying the result to make
    quantitative ensemble-size decisions.

    \item Use an LLM ensemble to navigate an ambiguous biomedical
    data dictionary, verify outputs against the authoritative
    database schema, and produce annotated executable code
    that implements a defensible query with all disambiguations
    documented.

    \item Evaluate the temporal inconsistency of racial and ethnic
    categories in the NCHS 1969--1986 dataset and produce a written
    race-variable handling strategy grounded in both the social
    construct framework and the specific coding history of the data.

    \item Produce a structured methodological consensus memo
    evaluating competing projection approaches, with
    FOO-calibrated confidence levels, a synthesis recommendation,
    and a retrospective comparison to the analytic plan developed
    in Unit~3.
\end{enumerate}

\subsection{Assessment Instrument}

\textbf{Item 3 --- Verified Data Query: Texas Birth Count
(Block~B deliverable):}
Annotated, executable code implementing the correct Texas birth
count filter. The annotation must explain: (a)~each filter field
used and why, (b)~each source of ambiguity in the data dictionary
and how it was resolved, (c)~the count produced, and (d)~at least
one invalidation from the ensemble interaction, logged in the
governance protocol format.

\textbf{Item 4 --- Invalidation Log (cumulative across Blocks~B,
C, and~D, and updated through Section~\ref{sec:7.3}):}
A structured log of all LLM invalidations encountered during
Unit~7. Each entry must record: invalidation type (factual,
logical, normative, or structural), a one-sentence description
of the failure, the resolution applied, and the section in which
it occurred. The log must contain at least six entries across
the full unit (minimum two from Section~\ref{sec:7.2}). If a
particular invalidation type does not appear in a team's
experience, this must be explicitly documented with a note on
what would have caused it to appear.

\textbf{Item 5 --- Race Variable Handling Strategy (Block~C
deliverable):}
A written strategy, two to three paragraphs, specifying: (a)~which
race-related variable(s) will be used in the projection analysis
and why, (b)~which temporal comparisons are valid within the
1969--1986 data and which are not, (c)~how the team's LLM queries
will be structured to mitigate normative invalidation on racial
disparity claims, and (d)~the updated race variable handling clause
for the governance protocol.

\textbf{Item 6 --- Methodological Consensus Memo (Block~D
deliverable):}
A structured memo containing: (a)~a comparison table with one
row per projection approach, recording ensemble consensus level
(High/Medium/Low), key assumptions, major risks, and divergence
points; (b)~the team's synthesis recommendation with justification;
(c)~the FOO analysis, which approach generated the most divergent
ensemble responses and what epistemic difficulty that signals;
and (d)~a retrospective comparison to the analytic plan designed
in Section~\ref{sec:3.1}, noting whether the consensus confirms,
complicates, or challenges the earlier design choices.

%================================================================
\section{AI Agents for Grant Review and Manuscript Preparation}
\label{sec:7.3}
%================================================================

\textbf{Time: 2 hours (Instructional: 20 minutes, Activity Work:
100 minutes)}

This session transfers the technical skills developed in
Section~\ref{sec:7.2} to the most career-relevant applications
of LLM-assisted research: grant review and manuscript preparation.
Students work with canonical NIH grant proposals that include
their actual panel summary statements, providing ground-truth
verification for ensemble outputs. Students also prepare a
Specific Aims page for a research proposal using Jackson Heart
Study (JHS) data, applying the race-consciousness framework from
Section~\ref{sec:7.1} to the framing of health disparity research.

The governance protocol from Section~\ref{sec:7.1} remains active.
The special confidentiality concern of Section~\ref{sec:7.3}A is
emphasized: federal agencies prohibit uploading unpublished
proposal text to unapproved AI systems for grant review purposes.
The proposals used here are canonical public examples; students
are explicitly warned that this constraint applies to any future
peer review panel work. The invalidation log from
Section~\ref{sec:7.2} is updated throughout.

\subsection*{Block A --- NIH Grant Review with Ground-Truth
Comparison (60 minutes)}

Each team selects one canonical NIH grant proposal from the
provided collection. These proposals are publicly available
examples accompanied by their actual summary statements from
review panels.

\textbf{Prompt design (10 min).} Teams design separate prompts
for each of the five NIH review criteria: Significance,
Investigator(s), Innovation, Approach, and Environment. Each
criterion requires different domain reasoning; a single prompt
for all five produces superficial output. Teams apply
Eq.~(\ref{eq:nmin}) per criterion, recognizing that Approach
typically warrants a higher agent count than Environment.

\textbf{Ensemble review (25 min).} Teams run the consensus
pipeline across all five criteria. The governance protocol is
active; every ensemble interaction is logged.

\textbf{Comparison analysis (20 min).} Teams compare the
ensemble's consensus review to the actual panel summary
statement across all five criteria, recording: consensus score,
panel score, strengths identified by each, divergence type
(missing context, conservative bias, detected concern missed
by panel, domain-norm gap), and agreement level.

\textbf{Critical synthesis (5 min).} Teams identify three
empirical categories from their comparison data: what the
ensemble reliably detected, what it missed systematically, and
what varied case by case. This analysis informs a critical
reflection on appropriate LLM use in grant preparation.

\subsection*{Block B --- JHS Specific Aims: Ensemble-Assisted
Drafting (60 minutes)}

The coordinator for JHS activities provides: a specific JHS
research question, the relevant JHS variables, and two to three
published JHS papers as grounding context.

\textbf{Prompt construction (10 min).} Teams design a structured
prompt that injects JHS context as a system prompt. The task
framing is explicit: not ``write my Specific Aims'' but rather
``given this JHS research question, this prior work, and these
three candidate aims, evaluate which aim structure is most
scientifically defensible and propose the structure with highest
consensus confidence, flagging any claims about JHS data or
methods that require independent verification against the
provided literature.''

\textbf{Ensemble drafting (25 min).} Teams run the FOO pipeline
with JHS context. Every LLM claim about JHS data, published
results, or methods is flagged for verification against the
provided papers and logged in the invalidation log.

\textbf{Human revision and race/ethnicity review (20 min).}
Teams revise the LLM draft, correcting all logged invalidations.
A mandatory step: teams explicitly evaluate whether the drafted
Specific Aims uses race and ethnicity appropriately as social
determinants of health rather than as biological variables, and
whether the framing aligns with current NIH policy on race
and ethnicity in grant applications. Each revision is documented
in the invalidation log.

\textbf{Disclosure statement finalization (5 min).} The draft
disclosure statement from Section~\ref{sec:7.1} is revised to
reflect the actual scope and nature of AI assistance across all
three sections. Teams select the appropriate disclosure level
(minimal, standard, or detailed) and produce a finalized
statement suitable for inclusion in a peer-reviewed publication.

\subsection{Learning Objectives}

\begin{enumerate}
    \item Design criterion-specific prompts for LLM-assisted NIH
    grant evaluation; compare ensemble outputs to actual expert
    panel judgments; and characterize empirical patterns of
    agreement, divergence, and systematic miss.

    \item Construct an ensemble-assisted Specific Aims draft for
    a JHS research proposal, verifying all LLM claims against
    published JHS literature and applying the race-consciousness
    framework to ensure appropriate framing of health disparity
    research.

    \item Produce a finalized disclosure statement appropriate for
    inclusion in a peer-reviewed publication, accurately reflecting
    the scope of AI assistance and maintaining full human
    accountability for all outputs.
\end{enumerate}

\subsection{Assessment Instrument}

\textbf{Item 7 --- NIH Grant Review Analysis (Block~A
deliverable):}
A structured comparison document containing: (a)~the comparison
table across all five NIH criteria with scores, identified
strengths, divergence types, and agreement levels; (b)~three
empirical categories (reliable detections, systematic misses,
and variable outcomes) derived from the comparison data with
at least one concrete example per category; and (c)~a one-paragraph
critical reflection on appropriate LLM use in grant preparation,
grounded in the evidence from the comparison rather than prior
assumptions.

\textbf{Item 8 --- JHS Specific Aims Draft (Block~B deliverable):}
A draft Specific Aims page produced through the ensemble pipeline
and revised by the team, with: (a)~all LLM invalidations
documented in the invalidation log with verification sources
cited; (b)~a written race/ethnicity review note confirming that
race and ethnicity are framed as social determinants of health
and that framing is consistent with NIH policy; and (c)~a brief
annotation for each aim identifying which elements are
human-authored and which were substantially shaped by ensemble
output, consistent with the governance protocol.

\textbf{Item 9 --- Final Disclosure Statement (Block~B
deliverable, unit-closing):}
A single, finalized disclosure statement appropriate for inclusion
in a methods section or acknowledgments section of a
peer-reviewed publication. The statement must: name the tools and
approximate versions used; specify the tasks for which AI
assistance was employed across the full unit; describe the
verification procedures applied; and include an explicit
accountability statement affirming that all claims, analyses,
and conclusions were reviewed and affirmed by the team, who
take full responsibility for the work. The statement must also
specify whether and how the LLM ensemble was used to review or
revise the predictive model report produced in Unit~4, including
a description of the verification procedure applied to ensemble
critiques of that report. The statement must be consistent with
the completed invalidation log.

% Required packages confirmed in DAIR3Config.tex:
% \usepackage{longtable}
% \usepackage{booktabs}
% \usepackage{ragged2e}
% \usepackage{array}
% \usepackage{pdflscape}

\begin{landscape}

    \section{Evaluation Rubric}

    {\scriptsize

    \begin{longtable}{%
        >{\RaggedRight\arraybackslash}p{2.5cm}   % Component
        >{\RaggedRight\arraybackslash}p{4.5cm}   % Excellent
        >{\RaggedRight\arraybackslash}p{4.2cm}   % Good
        >{\RaggedRight\arraybackslash}p{4.2cm}   % Satisfactory
        >{\RaggedRight\arraybackslash}p{4.5cm}}  % Needs Improvement

    \toprule
    \textbf{Component} &
    \textbf{Excellent (90--100\%)} &
    \textbf{Good (80--89\%)} &
    \textbf{Satisfactory (70--79\%)} &
    \textbf{Needs Improvement ($<$70\%)} \\
    \midrule
    \endfirsthead

    \toprule
    \textbf{Component} &
    \textbf{Excellent (90--100\%)} &
    \textbf{Good (80--89\%)} &
    \textbf{Satisfactory (70--79\%)} &
    \textbf{Needs Improvement ($<$70\%)} \\
    \midrule
    \endhead

    \midrule
    \multicolumn{5}{r}{\textit{Continued on next page}} \\
    \endfoot

    \bottomrule
    \endlastfoot

    % -------------------------------------------------------
    % Row 1: Invalidation Taxonomy Application (10%)
    % Maps to Assessment Item 1 (Section 7.1)
    % -------------------------------------------------------
    \textbf{1.\ Invalidation Taxonomy Application}
    \newline\textbf{(10\%)}
    &
    Correctly identifies all four invalidation types from the
    presented NCHS scenarios; provides precise, specific
    justification for each classification that demonstrates deep
    understanding of why the type applies and not another;
    proposes a concrete, type-appropriate verification method for
    each scenario; connects normative invalidation to the
    race-consciousness framework.
    &
    Correctly identifies most invalidation types; justifications
    are mostly accurate with minor gaps in specificity;
    verification methods proposed are appropriate with minor
    omissions; some connection to normative/race-consciousness
    dimension.
    &
    Identifies some invalidation types correctly but confuses
    at least one; justifications present but imprecise or
    incomplete; verification methods generic rather than
    type-specific; race-consciousness connection absent or
    superficial.
    &
    Fails to correctly identify multiple invalidation types or
    conflates types; justifications missing or incorrect;
    no meaningful verification strategies proposed; no engagement
    with normative dimension.
    \\
    \midrule

    % -------------------------------------------------------
    % Row 2: Research AI Governance Protocol (15%)
    % Maps to Assessment Item 2 (Section 7.1)
    % -------------------------------------------------------
    \textbf{2.\ Research AI Governance Protocol}
    \newline\textbf{(15\%)}
    &
    All five required components are present, complete, and
    specific to the NCHS/vital statistics context; provenance
    logging convention is operationalizable as written; race
    variable handling clause addresses coding history, temporal
    comparisons, and normative claim flagging; escalation
    thresholds are concrete and verifiable; disclosure statement
    follows the prescribed format and is appropriately specific.
    Protocol was demonstrably active and updated through
    Sections~7.2 and~7.3.
    &
    All five components present; most are specific and
    operationalizable with minor gaps; race clause addresses
    most key concerns; some evidence of protocol use in
    subsequent sections.
    &
    Most components present but one or more are generic or
    incomplete; race clause present but does not address
    temporal inconsistency or normative flagging substantively;
    limited evidence of protocol use beyond Section~7.1.
    &
    One or more required components missing; present components
    are too generic to operationalize; race clause absent
    or addresses only surface concerns; protocol not
    demonstrably used in subsequent activities.
    \\
    \midrule

    % -------------------------------------------------------
    % Row 3: Verified Data Query: Texas Birth Count (10%)
    % Maps to Assessment Item 3 (Section 7.2 Block B)
    % -------------------------------------------------------
    \textbf{3.\ Verified Data Query: Texas Birth Count}
    \newline\textbf{(10\%)}
    &
    Code is executable and produces the correct count; all
    ambiguities in the data dictionary are identified and
    resolved with explicit justification; annotations clearly
    explain each filter field and each disambiguation decision;
    at least one ensemble invalidation is documented in
    governance-protocol format; code structure is compatible
    with Unit~5 reproducible pipeline conventions.
    &
    Code is executable and produces a defensible count with
    minor filter issues; most ambiguities identified and
    resolved; annotations present for major decisions; at
    least one invalidation documented.
    &
    Code is present but contains a filter error or produces
    an uncorrected count; some ambiguities identified but
    not all resolved; annotations incomplete; invalidation
    documentation minimal.
    &
    Code is absent, non-executable, or implements an incorrect
    filter without acknowledgment; ambiguities in the data
    dictionary not identified; no invalidation documented.
    \\
    \midrule

    % -------------------------------------------------------
    % Row 4: Unit 4 Model Report Critique (10%)
    % Maps to Section 7.2 Block B model report critique sub-task
    % -------------------------------------------------------
    \textbf{4.\ Unit~4 Model Report Critique}
    \newline\textbf{(10\%)}
    &
    Ensemble critique is elicited with a well-designed prompt that
    directs attention to reproducibility, internal consistency, and
    data dictionary alignment; all flagged claims are classified
    by invalidation type with specific justification; each
    factual or structural invalidation is verified against
    Unit~4 code before acceptance; at least one false-negative
    or false-positive ensemble output is identified and
    documented in governance-protocol format; the invalidation
    log is updated consistently.
    &
    Critique prompt is adequate; most flagged claims classified
    by type with reasonable justification; verification performed
    for most invalidations; invalidation log updated.
    &
    Critique prompt is generic or insufficiently specific to
    the Unit~4 report; classification of flagged claims
    by type is partial or inconsistent; verification
    performed for some invalidations; log updated
    incompletely.
    &
    No critique prompt submitted, or critique not connected
    to the Unit~4 report; no classification by invalidation
    type; no verification against Unit~4 code; invalidation
    log not updated.
    \\
    \midrule

    % -------------------------------------------------------
    % Row 5: Invalidation Log (10%)
    % Maps to Assessment Item 4 (Sections 7.2 and 7.3)
    % -------------------------------------------------------
    \textbf{5.\ Invalidation Log}
    \newline\textbf{(10\%)}
    &
    Log contains at least six entries across the full unit with
    at least two from Section~7.2; all four invalidation types
    are represented or absent types are explicitly accounted for
    with a note on what would cause each to appear; every entry
    includes type, description, resolution, and section; log is
    consistent with the governance protocol format and with the
    disclosure statement.
    &
    Log contains at least six entries; most types represented
    or accounted for; most entries complete; minor inconsistency
    with governance protocol format.
    &
    Log contains fewer than six entries or multiple entries are
    incomplete; at least two invalidation types present;
    some entries lack type, description, or resolution.
    &
    Log contains fewer than three entries; entries are incomplete
    or mis-classified; absent types not accounted for; log
    is not consistent with governance protocol.
    \\
    \midrule

    % -------------------------------------------------------
    % Row 6: Race Variable Handling Strategy (10%)
    % Maps to Assessment Item 5 (Section 7.2 Block C)
    % -------------------------------------------------------
    \textbf{6.\ Race Variable Handling Strategy}
    \newline\textbf{(10\%)}
    &
    Strategy identifies a specific variable or combination
    with clear, data-dictionary-grounded justification;
    specifies which temporal comparisons are valid and which
    are not, with reference to the documented coding changes;
    concrete LLM query structure for mitigating normative
    invalidation is articulated; governance protocol clause
    is updated and consistent with the strategy.
    Strategy reflects understanding that race is a social
    construct whose reification has health consequences.
    &
    Variable selection justified with reasonable grounding;
    temporal validity addressed with minor gaps; normative
    invalidation mitigation described but not fully
    operationalized; protocol clause updated.
    &
    Variable selected but justification is weak or missing;
    temporal validity mentioned but not analyzed; normative
    invalidation acknowledged without concrete mitigation
    strategy; protocol clause minimally updated.
    &
    No specific variable identified or justification absent;
    temporal coding inconsistency not addressed; no normative
    invalidation mitigation; race treated as a natural
    biological category without acknowledgment of the social
    construct framework.
    \\
    \midrule

    % -------------------------------------------------------
    % Row 7: Methodological Consensus Memo (15%)
    % Maps to Assessment Item 6 (Section 7.2 Block D)
    % -------------------------------------------------------
    \textbf{7.\ Methodological Consensus Memo}
    \newline\textbf{(15\%)}
    &
    Comparison table is complete for all three approaches with
    accurate, specific entries for consensus level, assumptions,
    risks, and divergence points; synthesis recommendation is
    well-reasoned and explicitly connected to the data challenge
    context; FOO analysis identifies the most epistemically
    difficult approach and explains why divergence signals this;
    retrospective comparison to the Unit~3 analytic plan is
    substantive and either confirms or productively complicates
    the earlier design.
    &
    Table complete with mostly accurate entries; synthesis
    recommendation reasonable with adequate justification; FOO
    analysis present but lacks depth; retrospective comparison
    to Unit~3 present with minor gaps.
    &
    Table incomplete or contains inaccuracies for one approach;
    synthesis recommendation present but weakly justified;
    FOO analysis minimal; retrospective comparison to Unit~3
    superficial or absent.
    &
    Table missing one or more approaches or largely inaccurate;
    no synthesis recommendation or recommendation is unjustified;
    FOO analysis absent; no retrospective connection to Unit~3.
    \\
    \midrule

    % -------------------------------------------------------
    % Row 8: NIH Grant Review Analysis (10%)
    % Maps to Assessment Item 7 (Section 7.3 Block A)
    % -------------------------------------------------------
    \textbf{8.\ NIH Grant Review Analysis}
    \newline\textbf{(10\%)}
    &
    Comparison table covers all five criteria with accurate,
    evidence-grounded entries; three empirical categories
    (reliable detections, systematic misses, variable outcomes)
    are clearly distinguished with at least one concrete example
    each; critical reflection is grounded in comparison evidence
    rather than prior assumptions and reaches a nuanced,
    defensible conclusion about appropriate LLM use in grant
    preparation.
    &
    Table covers all five criteria with mostly accurate entries;
    three categories identified with examples, some less
    concrete; reflection grounded in evidence with minor gaps
    in nuance.
    &
    Table covers most criteria; categories identified but
    blurred or with weak examples; reflection present but
    relies on assertion more than evidence.
    &
    Table incomplete or inaccurate for multiple criteria;
    categories not meaningfully distinguished; reflection
    absent or disconnected from comparison evidence.
    \\
    \midrule

    % -------------------------------------------------------
    % Row 9: JHS Specific Aims Draft (10%)
    % Maps to Assessment Item 8 (Section 7.3 Block B)
    % -------------------------------------------------------
    \textbf{9.\ JHS Specific Aims Draft}
    \newline\textbf{(10\%)}
    &
    Draft is complete (opening paragraph, three aims, closing
    paragraph); all LLM invalidations are logged with
    verification sources; race/ethnicity review note confirms
    appropriate social-determinant framing and NIH policy
    alignment with specific evidence from the draft;
    annotation accurately distinguishes human-authored from
    ensemble-shaped elements; draft is scientifically coherent
    and grounded in the provided JHS literature.
    &
    Draft is mostly complete; most invalidations logged with
    sources; race/ethnicity framing mostly appropriate with
    minor lapses; annotation present with minor gaps; draft
    is scientifically reasonable.
    &
    Draft has structural gaps (missing aim or section);
    invalidation log incomplete or sources missing; race/
    ethnicity framing partially addressed without explicit
    review note; annotation minimal.
    &
    Draft is incomplete or incoherent; invalidation log
    absent or has fewer than two entries; no race/ethnicity
    review performed; no annotation distinguishing human
    from ensemble contributions.
    \\
    \midrule

    % -------------------------------------------------------
    % Row 10: Final Disclosure Statement (5%)
    % Maps to Assessment Item 9 (Section 7.3 Block B, closing)
    % -------------------------------------------------------
    \textbf{10.\ Final Disclosure Statement}
    \newline\textbf{(5\%)}
    &
    Statement names tools and approximate versions; accurately
    specifies the tasks for which AI assistance was employed
    across all three sections; describes the verification
    procedures applied in terms consistent with the
    invalidation log; includes an explicit accountability
    statement affirming full human responsibility; is appropriate
    for inclusion in a peer-reviewed methods or acknowledgments
    section without revision.
    &
    Statement names tools and tasks with minor omissions;
    verification procedures described at a general level;
    accountability statement present; suitable for publication
    with minor editing.
    &
    Statement present but omits tools or tasks; verification
    described only vaguely; accountability statement present
    but formulaic; would require substantial editing for
    publication.
    &
    Statement absent, or present but does not name tools,
    tasks, or verification; no accountability statement;
    not suitable for publication in any form.
    \\

    \end{longtable}

    } % end \footnotesize

\end{landscape}
